\documentclass[11pt,a4paper]{article}

% ============================================================
%  Preámbulo común --- Tarea LEC06 CC451
% ============================================================

% --- Idioma y codificación ---
\usepackage[utf8]{inputenc}
\usepackage[T1]{fontenc}
\usepackage[spanish,es-tabla]{babel}
\usepackage{lmodern}
\usepackage{microtype}

% --- Geometría ---
\usepackage[a4paper,margin=2.4cm]{geometry}
\usepackage{setspace}
\onehalfspacing

% --- Tipografía y enlaces ---
\usepackage{xcolor}
\definecolor{azulUNI}{RGB}{0,51,102}
\definecolor{grisCajas}{RGB}{240,240,240}
\definecolor{grisBorde}{RGB}{200,200,200}

\usepackage[colorlinks=true,
            linkcolor=azulUNI,
            citecolor=azulUNI,
            urlcolor=azulUNI]{hyperref}

% --- Tablas y listas ---
\usepackage{booktabs}
\usepackage{tabularx}
\usepackage{longtable}
\usepackage{array}
\usepackage{multirow}
\usepackage{enumitem}

% --- Figuras ---
\usepackage{graphicx}
\usepackage{float}
\graphicspath{{figuras/}}

% --- Diagrama de Gantt ---
\usepackage{pgfgantt}
\setlist{nosep,leftmargin=1.4em}

% --- Encabezados de sección ---
\usepackage{titlesec}
\titleformat{\section}
  {\color{azulUNI}\Large\bfseries}{\thesection}{0.6em}{}
\titleformat{\subsection}
  {\color{azulUNI}\large\bfseries}{\thesubsection}{0.5em}{}
\titleformat{\subsubsection}
  {\bfseries}{\thesubsubsection}{0.4em}{}

% --- Cabecera y pie ---
\usepackage{fancyhdr}
\pagestyle{fancy}
\fancyhf{}
\fancyhead[L]{\small CC451 --- IHC --- LEC06}
\fancyhead[R]{\small UNI \quad 2026-I}
\fancyfoot[C]{\thepage}
\renewcommand{\headrulewidth}{0.3pt}

% --- Cajas de aviso y código ---
\usepackage{tcolorbox}
\tcbuselibrary{breakable,skins}

\newtcolorbox{cajaPrompt}[1][]{%
  enhanced,
  colback=grisCajas,
  colframe=azulUNI,
  boxrule=0.5pt,
  arc=2pt,
  fonttitle=\bfseries\color{white},
  coltitle=white,
  attach boxed title to top left={xshift=4mm,yshift=-2mm},
  boxed title style={colback=azulUNI,sharp corners},
  breakable,
  before skip=10pt, after skip=10pt,
  #1
}

\newtcolorbox{cajaCriterio}[1][]{%
  enhanced,
  colback=white,
  colframe=grisBorde,
  boxrule=0.4pt,
  arc=1pt,
  breakable,
  before skip=6pt, after skip=10pt,
  #1
}

% --- Listings (para mostrar prompts en monoespaciado si conviene) ---
\usepackage{listings}
\lstset{
  basicstyle=\ttfamily\footnotesize,
  breaklines=true,
  breakatwhitespace=true,
  frame=single,
  rulecolor=\color{grisBorde},
  backgroundcolor=\color{grisCajas},
  columns=fullflexible,
}

% --- Comandos cortos ---
\newcommand{\rf}[1]{\textbf{RF-#1}}
\newcommand{\rnf}[1]{\textbf{RNF-#1}}
\newcommand{\cu}[1]{\textbf{CU-#1}}
\newcommand{\sk}[1]{\textbf{S#1}}

% Tabularx column type (J = justified)
\newcolumntype{Y}{>{\raggedright\arraybackslash}X}


\title{%
  \vspace{-1.5cm}
  \large CC451 --- Interacción Humano Computadora \\[0.3cm]
  \Large Tarea LEC06 --- Definición Formal del Proyecto \\[0.5cm]
  \LARGE\bfseries Simulador de Entrevistas Laborales Adaptativo \\[0.3cm]
  \large \textit{Reclutador LLM con análisis multimodal de voz y lenguaje corporal}
}

\author{%
  Aaron Davila Santos \and
  Max Serrano Arostegui \and
  Walter Poma Navarro
}

\date{Universidad Nacional de Ingeniería \\
      Profesor: Ciro Núñez Iturri \\
      Ciclo 2026-I \quad --- \quad 29 de abril de 2026}

\begin{document}

\maketitle
\thispagestyle{empty}

\begin{abstract}
\noindent Este documento constituye la definición formal (PC02) del proyecto del curso CC451.
Se presentan: (i) el tema y objetivo principal, (ii) el listado completo de
\textit{stakeholders} incluyendo usuarios con discapacidad, (iii) los requerimientos
funcionales y no funcionales asociados a cada uno, (iv) el cumplimiento de los
módulos UX/UI obligatorios y opcionales exigidos por el enunciado, y (v) las
especificaciones detalladas de los casos de uso. La elaboración del documento
utilizó dos proveedores de IA generativa (Google Gemini y Anthropic Claude,
ambos en versión web) cuyos prompts optimizados y criterios de selección de
resultados se reportan al final.
\end{abstract}

\tableofcontents
\newpage

\section{Tema y objetivo del proyecto}
\label{sec:tema}

\subsection{Título}
\textbf{Simulador de Entrevistas Laborales Adaptativo}: un reclutador basado en
LLM que adapta sus preguntas a la industria y nivel del candidato, analizando
en tiempo real la voz (fluidez, muletillas, pausas, tono) y el lenguaje
corporal por webcam (contacto visual, postura, gestualidad), y que habilita
\textit{mock interviews} colaborativas entre estudiantes con roles
intercambiables.

\subsection{Objetivo principal}
Diseñar y construir una aplicación web/PWA que permita a estudiantes y
recién egresados practicar entrevistas laborales en cualquier momento,
recibiendo retroalimentación objetiva, instantánea y multimodal sobre su
desempeño verbal y no verbal, así como un plan de mejora personalizado y la
opción de validación entre pares.

\subsection{Objetivos específicos}
\begin{enumerate}
  \item Implementar un agente entrevistador conversacional por voz capaz de
        adaptar dinámicamente la dificultad y el dominio (industria, rol,
        nivel) según el perfil del candidato.
  \item Calcular y mostrar en tiempo real métricas de comunicación verbal
        (palabras-muletilla por minuto, ritmo, claridad) y no verbal
        (contacto visual, postura, gestualidad) sin enviar video crudo a
        servicios externos.
  \item Generar un plan de mejora individualizado tras cada sesión, basado
        en debilidades recurrentes detectadas a lo largo del historial
        longitudinal del usuario.
  \item Habilitar sesiones \textit{peer-to-peer} entre estudiantes mediante
        WebRTC, con panel de observador, comentarios anclados a
        \textit{timestamps} e intercambio de roles.
  \item Aplicar elementos de gamificación (rangos por competencia, ligas
        temáticas, \textit{badges}) que sostengan la motivación intrínseca
        sin recurrir a \textit{dark patterns}.
  \item Garantizar accesibilidad WCAG 2.2 nivel AA y rutas alternativas para
        usuarios con discapacidad visual, auditiva, motora o cognitiva.
\end{enumerate}

\subsection{Problema real que resuelve}
Los estudiantes y recién egresados enfrentan dos barreras al prepararse para
entrevistas técnicas:

\begin{itemize}
  \item \textbf{Falta de retroalimentación objetiva}: amigos y familiares no
        pueden evaluar lenguaje corporal y fluidez de manera consistente, y
        el sesgo afectivo distorsiona su juicio.
  \item \textbf{Costo y escasez de coaches profesionales}: el coaching uno
        a uno es caro, requiere agenda y suele ser inaccesible para
        estudiantes universitarios.
\end{itemize}

\noindent Una herramienta multimodal con IA permite practicar
$24\times 7$, sin exposición inicial a juicio humano, y luego con
validación opcional entre pares.

\subsection{Diferenciación frente a propuestas similares}
A diferencia de aplicaciones para entrenamiento de habilidades sociales
orientadas al bienestar emocional, esta propuesta:

\begin{itemize}
  \item Se enfoca en empleabilidad técnica (entrevistas reales de trabajo).
  \item Usa análisis audiovisual multimodal en tiempo real (voz $+$ video),
        no solo diálogo textual.
  \item Incorpora una dimensión comunitaria y competitiva (peer mock,
        ligas, \textit{ranking}).
  \item Su métrica objetivo es desempeño profesional, no reducción de
        síntomas clínicos.
\end{itemize}

\section{Stakeholders del proyecto}
\label{sec:stakeholders}

Esta sección consolida la lista de \textit{stakeholders} a partir del
prompt 01 (Gemini), enriquecida por el equipo con actores que el modelo
no incluyó (\textit{peer} entrevistador como rol activo,
reclutador/coach como consultor, administrador del sistema como rol
diferenciado del equipo de desarrollo, y comunidad académica de IHC).
La numeración resultante \sk{1}--\sk{17} es la \textbf{autoritativa} en
todo el documento y se referencia en las secciones de requerimientos
(\ref{sec:requerimientos}) y casos de uso (\ref{sec:casos}).

Clasificación: \textbf{P} = primario (interactúa directamente con la
aplicación), \textbf{S} = secundario (interactúa de forma indirecta),
\textbf{T} = terciario (afectado pero no interactúa).

\subsection{Listado completo}

\begin{longtable}{p{1.0cm} p{0.6cm} p{4.4cm} p{8.0cm}}
\toprule
\textbf{ID} & \textbf{Tipo} & \textbf{Stakeholder} & \textbf{Interés / rol respecto al sistema} \\
\midrule
\endhead

\sk{1}  & P & Estudiante de ingeniería en búsqueda de prácticas
preprofesionales (UNI) &
Interactúa con el sistema para entrenar habilidades blandas y técnicas
específicas del perfil de egreso. Su interés es obtener feedback
accionable para reducir la ansiedad y mejorar su desempeño en procesos
reales. \\
\addlinespace

\sk{2}  & P & Recién egresado en búsqueda de su primer empleo a tiempo
completo &
Utiliza la plataforma para simular entrevistas de mayor complejidad
técnica y estratégica según el mercado laboral actual. Busca validar
su capacidad de comunicación efectiva y el manejo de lenguaje corporal
frente a un evaluador basado en IA. \\
\addlinespace

\sk{3}  & P & Usuario en rol de \textit{peer} entrevistador &
El mismo estudiante (\sk{1} o \sk{2}) en rol invertido durante
\textit{mock interviews} colaborativas: formula preguntas, ancla
comentarios, califica. \\
\addlinespace

\sk{4}  & P & Usuario con discapacidad visual (baja visión o ceguera) &
Accede a la interfaz mediante lectores de pantalla y comandos de voz,
requiriendo cumplimiento de WCAG 2.2 nivel AA. Su interés es el
entrenamiento de la voz y el ritmo sin barreras de accesibilidad
visual. \\

\sk{5}  & P & Usuario con discapacidad auditiva (hipoacusia o sordera) &
Interactúa mediante subtitulado en tiempo real y señales visuales que
compensan la retroalimentación sonora. Depende de una interfaz que
traduzca el tono y ritmo del entrevistador en indicadores claros. \\

\sk{6}  & P & Usuario con discapacidad motora (limitación funcional en
extremidades superiores) &
Utiliza la aplicación mediante tecnologías de asistencia como
pulsadores o control por voz. Requiere que la interfaz no dependa de
interacciones físicas complejas o tiempos de respuesta táctiles
exigentes. \\

\sk{7}  & P & Usuario con condición de neurodivergencia (TEA, TDAH,
dislexia) &
Emplea el simulador para practicar el contacto visual y la
interpretación de señales sociales en un entorno controlado y
predecible. Su interés principal es la reducción de la sobrecarga
cognitiva mediante una interfaz clara y la repetición ilimitada de
sesiones. \\
\addlinespace

\sk{8}  & S & Reclutador o profesional de RRHH (consultor del proyecto) &
Aporta criterios reales de evaluación de entrevistas; revisa la
calibración de las preguntas y rúbricas generadas por el LLM. \\
\addlinespace

\sk{9}  & S & Coach de carrera o mentor &
Recomienda la herramienta, revisa el plan de mejora del candidato y
puede dejar comentarios asincrónicos en sesiones grabadas. \\
\addlinespace

\sk{10} & S & Oficina de Empleabilidad de la Universidad Nacional de
Ingeniería &
Monitorea reportes agregados y anonimizados para identificar brechas
de habilidades comunes en los estudiantes. Utiliza estos datos para
ajustar talleres de inserción laboral y mejorar los indicadores de
empleabilidad de la institución. \\
\addlinespace

\sk{11} & S & Equipo de desarrollo y mantenimiento del proyecto &
Responsable de la implementación técnica, mantenimiento de los modelos
de MediaPipe y la integración de las API multimodales. Asegura
estabilidad operativa del backend y persistencia de los planes de
mejora personalizados. \\
\addlinespace

\sk{12} & S & Administrador del sistema &
Gestiona usuarios, ligas, contenido de industrias y moderación del
\textit{feed} comunitario. Diferenciado del equipo de desarrollo
(\sk{11}) por tener acceso a operaciones en producción y no a
despliegues. \\
\addlinespace

\sk{13} & S & Docentes del curso CC451 (cátedra) &
Supervisan el uso del sistema como herramienta pedagógica para
evaluación de métricas de UX y accesibilidad. Validan que el proyecto
cumpla con los objetivos académicos y técnicos planteados en el ciclo
2026-I. \\
\addlinespace

\sk{14} & S & Administradores de infraestructura en la nube
(Render, Fly.io, Vercel) &
Proveen servicios de cómputo y red para el despliegue de la PWA y el
backend. Su interés es el consumo eficiente de recursos y el
cumplimiento de los acuerdos de nivel de servicio. \\
\addlinespace

\sk{15} & T & Empresas reclutadoras del sector tecnológico y de
ingeniería &
Beneficiarias indirectas: reciben candidatos mejor preparados. No
interactúan directamente con el software, pero el resultado del
entrenamiento impacta en la calidad de sus procesos de selección. \\
\addlinespace

\sk{16} & T & Autoridad Nacional de Protección de Datos Personales
(Perú, ANPDP) &
Vela por el cumplimiento de la Ley N\textsuperscript{\textordmasculine}~29733
respecto al tratamiento de datos biométricos y grabaciones procesadas
por el sistema. Su rol es regulatorio: asegurar que no se vulnere la
privacidad de los estudiantes durante el análisis de voz y video. \\
\addlinespace

\sk{17} & T & Familiares y entorno cercano del candidato &
Se ven influenciados por la mejora en las perspectivas económicas y el
éxito profesional del usuario tras el uso del simulador. Representan
el entorno de apoyo que se beneficia indirectamente de la inserción
laboral efectiva del candidato. \\

\bottomrule
\end{longtable}

\subsection{Mapa de prioridades}
Para el alcance del MVP se priorizan los \textit{stakeholders}
primarios (\sk{1}--\sk{7}), garantizando que las rutas de uso de los
cuatro perfiles con discapacidad (\sk{4}--\sk{7}) sean equivalentes en
funcionalidad a la del candidato general (\sk{1}). Los
\textit{stakeholders} secundarios y terciarios influyen en
requerimientos no funcionales (privacidad, mantenibilidad, costo de
operación, cumplimiento normativo) y se atienden en la
sección~\ref{sec:requerimientos}.

\subsection{Equivalencia con la numeración propuesta por Gemini}
La respuesta cruda del prompt 01 (versionada en
\texttt{prompts/respuestas/gemini\_01\_stakeholders.md}) usa una
numeración propia de 14 \textit{stakeholders}. La equivalencia con la
numeración autoritativa del informe es:

\begin{itemize}
  \item Gemini-S1 $=$ \sk{1}; Gemini-S2 $=$ \sk{2}.
  \item Gemini-S3, S4, S5, S6 $=$ \sk{4}, \sk{5}, \sk{6}, \sk{7}
        (perfiles con discapacidad).
  \item Gemini-S7 $=$ \sk{10} (Oficina de Empleabilidad UNI).
  \item Gemini-S8 $=$ \sk{13} (Docentes del curso).
  \item Gemini-S9 $=$ \sk{11} (equipo de desarrollo).
  \item Gemini-S10 $=$ \sk{14} (infraestructura cloud).
  \item Gemini-S11 $=$ \sk{15} (empresas reclutadoras).
  \item Gemini-S12 $=$ \sk{16} (autoridad de protección de datos).
  \item Gemini-S13 $=$ \textit{proveedor LLM externo}, no incluido como
        \textit{stakeholder} en el informe; aparece como dependencia
        técnica en arquitectura.
  \item Gemini-S14 $=$ \sk{17} (familiares).
\end{itemize}

\noindent El equipo agregó cuatro \textit{stakeholders} no propuestos
por Gemini: \sk{3} (peer entrevistador como rol activo), \sk{8}
(reclutador/RRHH como consultor del proyecto), \sk{9} (coach/mentor
con feedback asincrónico) y \sk{12} (administrador del sistema
diferenciado del equipo de desarrollo).

\section{Usuarios con discapacidad y funcionalidad asociada}
\label{sec:discapacidad}

Los \textit{stakeholders} \sk{4}--\sk{7} forman parte del usuario nuclear
del simulador. La estrategia es ofrecer \textbf{equivalencia funcional}:
toda capacidad del candidato general (\sk{1}) debe estar disponible para
ellos por una ruta accesible. Se trabaja sobre WCAG 2.2 nivel AA, con
verificación automática en CI mediante \texttt{axe-core} y
\texttt{Lighthouse}, y validación manual con representantes de cada
perfil. Esta sección integra la respuesta del proveedor Claude
(prompt 02), seleccionada por su rigor en la cita de criterios WCAG
concretos y por proponer métricas sustitutas que evitan penalizar
injustamente al usuario.

\subsection{Discapacidad visual (\sk{4})}

\subsubsection*{A. Necesidades específicas}
\begin{itemize}
  \item Acceso completo al contenido textual y gráfico de la sala
        virtual mediante lector de pantalla (NVDA, JAWS, VoiceOver),
        incluyendo el estado del aura del entrevistador y los
        indicadores de progreso de la sesión.
  \item Sustitución del canal visual del aura por un canal sonoro o
        textual equivalente, ya que los degradados luminosos son
        inaccesibles para usuarios ciegos y de bajo contraste
        discriminable para usuarios con baja visión.
  \item Control de tipografía, contraste y zoom del cliente sin que se
        rompan los flujos de entrevista, conservando la persistencia
        de configuración entre sesiones.
  \item Eliminación de la métrica de contacto visual del cómputo del
        puntaje global, sustituida por una métrica de orientación
        auditiva equivalente.
  \item Información clara y anticipada sobre el uso de la cámara:
        posibilidad de mantenerla apagada sin perder funcionalidades
        evaluables.
\end{itemize}

\subsubsection*{B. Funcionalidades verificables}
\begin{enumerate}
  \item Operación al 100\,\% por teclado en todos los flujos
        (login, configuración, entrevista, revisión de feedback, peer
        mock), conforme al criterio WCAG 2.2 2.1.1 \textit{Keyboard};
        verificable mediante recorrido E2E con Playwright sin uso de
        puntero.
  \item Todo elemento interactivo expone nombre, rol y estado
        accesibles vía ARIA, conforme al criterio WCAG 2.2 4.1.2
        \textit{Name, Role, Value}; verificable con \texttt{axe-core}
        en CI con cero violaciones de severidad
        \textit{serious}/\textit{critical}.
  \item El aura luminosa se traduce en un canal sonoro paralelo
        (\textit{earcons} no verbales con mapeo frecuencia-fluidez y
        panorámica-ritmo) que se activa con un \textit{toggle}
        persistido en el perfil; verificable con E2E que valida la
        emisión del \textit{earcon} ante eventos sintéticos.
  \item La métrica de contacto visual se reemplaza automáticamente por
        \emph{orientación de voz hacia el micrófono} cuando el usuario
        declara discapacidad visual o desactiva la cámara; verificable
        por telemetría del cambio de métrica activa y por test
        unitario del módulo de \textit{scoring}.
  \item El cliente cumple WCAG 2.2 1.4.3 \textit{Contrast (Minimum)}
        con ratio mínimo de 4.5:1 en texto y 3:1 en componentes, y
        WCAG 2.2 1.4.4 \textit{Resize Text} permitiendo zoom hasta
        200\,\% sin pérdida de contenido ni funcionalidad; verificable
        con Lighthouse Accessibility (puntaje mínimo 95) y pruebas
        visuales con zoom forzado.
  \item El feedback final se entrega también como reporte estructurado
        en HTML semántico y descargable en texto plano, navegable por
        encabezados y regiones; verificable con \texttt{axe-core} sobre
        el reporte y con E2E que confirma jerarquía h1--h2--h3 sin
        saltos.
\end{enumerate}

\subsection{Discapacidad auditiva (\sk{5})}

\subsubsection*{A. Necesidades específicas}
\begin{itemize}
  \item Equivalente textual en tiempo real para toda emisión de voz del
        entrevistador LLM, con latencia compatible con la cadencia
        conversacional de una entrevista.
  \item Posibilidad de responder en lengua de señas o por escrito sin
        que se degrade el análisis de competencias comunicativas.
  \item Sustitución de las métricas dependientes de prosodia (tono,
        pausas, muletillas) por métricas observables en el canal
        alternativo elegido por el usuario.
  \item Indicadores visuales redundantes para cualquier alerta o
        evento que en el flujo estándar se señaliza por audio.
  \item Compatibilidad con auriculares con bobina inductiva (T-coil) y
        control independiente del volumen del entrevistador respecto al
        volumen de notificaciones.
\end{itemize}

\subsubsection*{B. Funcionalidades verificables}
\begin{enumerate}
  \item Subtítulos sincronizados de las preguntas del entrevistador
        mediante TTS con \textit{captions} WebVTT, latencia mediana
        $\le 500\,\mathrm{ms}$ \texttt{(referencial)}, conforme a
        WCAG 2.2 1.2.4 \textit{Captions (Live)}; verificable por
        telemetría de latencia y E2E que valida la presencia y
        sincronía de pistas WebVTT.
  \item Modo \emph{respuesta escrita}: el turno se cierra mediante
        envío de texto en lugar de detección de fin de habla;
        verificable con E2E que completa una entrevista íntegra en
        este modo y telemetría que registra el modo activo por sesión.
  \item Cuando se selecciona discapacidad auditiva o el modo escrito,
        las métricas de fluidez verbal y muletillas se sustituyen por
        \emph{coherencia y densidad léxica} sobre la transcripción del
        usuario; verificable con test unitario del módulo de
        \textit{scoring}.
  \item Todos los eventos de audio (turno asignado, alerta de tiempo,
        cambio de fase, desconexión en peer mock) tienen contraparte
        visual persistente de al menos $3\,\mathrm{s}$
        \texttt{(referencial)}, conforme a WCAG 2.2 1.4.13
        \textit{Content on Hover or Focus}; verificable con E2E que
        dispara eventos sintéticos y valida la aparición en el DOM.
  \item El componente de subtítulos cumple WCAG 2.2 1.4.12
        \textit{Text Spacing} y permite ajuste de tamaño, color de
        fondo y opacidad sin recargar la sesión; verificable con
        \texttt{axe-core} y E2E de persistencia entre sesiones.
  \item La sala de peer mock incluye un canal de chat textual paralelo
        al canal de voz, con historial accesible durante toda la
        sesión; verificable con E2E multiusuario que confirma entrega
        bidireccional en menos de $1\,\mathrm{s}$ \texttt{(referencial)}.
\end{enumerate}

\subsection{Discapacidad motora (\sk{6})}

\subsubsection*{A. Necesidades específicas}
\begin{itemize}
  \item Operación del simulador sin requerir gestos finos de mouse,
        soportando teclado, conmutadores (\textit{switch access}) y
        comandos de voz.
  \item Tolerancia a tiempos de respuesta variables: ningún flujo
        crítico debe expirar por inactividad sin posibilidad de
        extensión.
  \item Eliminación o sustitución de las métricas de postura y
        gesticulación, que penalizarían a usuarios con movilidad
        reducida del tronco, brazos o cabeza.
  \item Encuadre flexible de la cámara: aceptación de planos no
        convencionales (lateral, sentado en silla de ruedas, decúbito)
        sin que la detección considere la imagen como inválida.
  \item Áreas de activación amplias y sin requerimiento de pulsaciones
        simultáneas o sostenidas.
\end{itemize}

\subsubsection*{B. Funcionalidades verificables}
\begin{enumerate}
  \item Toda acción ejecutable mediante un único punto de entrada
        (teclado, \textit{switch} o voz) sin combinaciones simultáneas
        obligatorias, conforme a WCAG 2.2 2.1.4
        \textit{Character Key Shortcuts} y 2.5.7
        \textit{Dragging Movements}; verificable con E2E
        \textit{keyboard-only} y otro \textit{voice-only}.
  \item Todos los objetivos táctiles tienen tamaño mínimo de
        $24\times 24$\,px CSS, conforme al criterio WCAG 2.2 2.5.8
        \textit{Target Size (Minimum)}; verificable con auditoría
        automatizada que mide el \textit{bounding box} de todos los
        elementos interactivos.
  \item Los temporizadores de respuesta del entrevistador son
        configurables (sin límite, estándar, $\times 2$, $\times 3$) y
        notifican antes de expirar con opción de prórroga, conforme a
        WCAG 2.2 2.2.1 \textit{Timing Adjustable}; verificable con E2E
        y telemetría del temporizador efectivo por sesión.
  \item Cuando el usuario declara discapacidad motora, las métricas de
        postura y gesticulación se reemplazan por
        \emph{estabilidad del encuadre facial} (presencia y centrado
        relativo del rostro) sin penalizar inclinaciones, asimetrías
        ni ausencia de movimiento corporal; verificable con test
        unitario del \textit{pipeline} MediaPipe y regresión sobre
        \textit{datasets} sintéticos de planos no convencionales.
  \item Capa de comandos de voz para todos los controles
        (\textit{``siguiente pregunta''}, \textit{``pausar''},
        \textit{``repetir''}), independiente del canal de respuesta
        sustantiva, conforme a WCAG 2.2 2.5.4 \textit{Motion
        Actuation}; verificable con E2E que dispara los comandos por
        API sintética y valida los efectos en el DOM.
  \item El sistema permite recalibrar el encuadre de la cámara
        declarando explícitamente la posición del usuario (de frente,
        lateral, sentado bajo, recostado) sin que la detección marque
        la imagen como no válida; verificable con test de integración
        que inyecta \textit{frames} de cada posición.
\end{enumerate}

\subsection{Discapacidad cognitiva y neurodivergencia (\sk{7})}

\subsubsection*{A. Necesidades específicas}
\begin{itemize}
  \item Reducción configurable de estímulos visuales y sonoros: el
        aura, las animaciones de fondo y los efectos de transición
        pueden saturar a usuarios con TEA o TDAH.
  \item Lenguaje claro en preguntas y feedback, con posibilidad de
        reformular en versión simplificada sin alterar el contenido
        evaluativo.
  \item Soporte tipográfico para usuarios con dislexia (fuentes
        específicas, espaciado aumentado, alineación a la izquierda
        sin justificado).
  \item Previsibilidad y control sobre la duración y estructura de la
        sesión, evitando cambios de fase sin aviso.
  \item Posibilidad de pausar la sesión sin penalización ni pérdida de
        progreso, en respuesta a sobrecarga sensorial o de atención.
\end{itemize}

\subsubsection*{B. Funcionalidades verificables}
\begin{enumerate}
  \item Modo de \emph{estimulación reducida} que desactiva el aura, las
        animaciones de fondo, los \textit{earcons} no esenciales y las
        transiciones, dejando únicamente indicadores estáticos de
        estado; verificable con E2E que confirma la ausencia de
        animaciones CSS activas (\texttt{prefers-reduced-motion}
        honrado), conforme a WCAG 2.2 2.3.3 \textit{Animation from
        Interactions}, y con auditoría Lighthouse.
  \item El usuario puede limitar el aura a un máximo de dos métricas
        de su elección entre las disponibles; verificable con E2E que
        selecciona la configuración y valida que el componente del
        aura solo renderiza los canales seleccionados.
  \item Todas las preguntas del entrevistador disponen de una versión
        en lenguaje claro (oraciones cortas, voz activa, una idea por
        oración) accesible mediante un control persistente; verificable
        con test unitario sobre el módulo de generación y métrica de
        legibilidad automatizada (índice INFLESZ o similar) con puntaje
        mínimo \texttt{(referencial)}.
  \item Control de tipografía con opciones validadas para dislexia
        (fuente con \textit{tracking} aumentado y espaciado interlínea
        $\ge 1.5$), conforme a WCAG 2.2 1.4.12 \textit{Text Spacing};
        verificable con E2E que aplica la configuración y valida
        \texttt{line-height}, \texttt{letter-spacing} y
        \texttt{word-spacing} en el DOM.
  \item Pausa y reanudación en cualquier punto sin pérdida del estado
        conversacional ni penalización en el puntaje, con estado
        conservado por al menos $24\,\mathrm{h}$ \texttt{(referencial)},
        conforme a WCAG 2.2 2.2.1 \textit{Timing Adjustable} y 2.2.6
        \textit{Timeouts}; verificable con E2E que pausa, cierra el
        navegador, reabre y valida la continuidad del contexto.
  \item Antes de iniciar cada fase, el simulador anuncia su duración
        estimada, su estructura y los criterios que se evaluarán, en
        formato textual persistente, conforme a la práctica de
        previsibilidad asociada a WCAG 2.2 3.2.5 \textit{Change on
        Request}; verificable con E2E que confirma la presencia del
        bloque de anuncio en cada transición de fase.
\end{enumerate}

\subsection{Verificación}

La validación combina cuatro herramientas y cuatro hitos por perfil:

\paragraph{Herramientas} \texttt{axe-core} integrado en CI con umbral
cero para violaciones \textit{serious} y \textit{critical};
\texttt{Lighthouse} con puntaje Accessibility mínimo de 95;
\texttt{Playwright} para E2E que cubren los flujos
\textit{keyboard-only}, modo escrito, comandos de voz y configuración
de tipografía; panel de telemetría propio que registra modo activo,
métricas sustitutas, latencia de subtítulos y temporizadores efectivos.

\paragraph{Hitos}
\begin{enumerate}
  \item Revisión heurística experta antes de la primera prueba con
        usuarios.
  \item Prueba moderada con un usuario representativo por perfil
        (\sk{4}--\sk{7}) sobre una entrevista completa.
  \item Peer mock con al menos un par mixto que incluya un usuario
        del perfil correspondiente.
  \item Auditoría final externa con conformidad WCAG 2.2 nivel AA
        documentada en informe firmado, previa a la entrega.
\end{enumerate}

\section{Requerimientos funcionales y no funcionales}
\label{sec:requerimientos}

Cada requerimiento se identifica con \rf{n} (funcional) o \rnf{n} (no
funcional), se asocia a uno o más \textit{stakeholders} y debe ser
\textbf{automatizable} por la aplicación (verificable mediante prueba
automática, métrica de telemetría o validación de la UI). La trazabilidad
con casos de uso aparece en la sección \ref{sec:casos}.

\subsection{Requerimientos funcionales}

\begin{longtable}{p{1.4cm} p{8.0cm} p{2.0cm} p{2.6cm}}
\toprule
\textbf{ID} & \textbf{Descripción} & \textbf{Stakeholders} & \textbf{Verificación automática} \\
\midrule
\endhead

\rf{01} & Iniciar una sesión de entrevista seleccionando industria, rol
y nivel objetivo. & \sk{1}--\sk{7} & Test E2E de creación de sesión \\
\addlinespace
\rf{02} & Sostener un diálogo bidireccional por voz con un avatar
entrevistador basado en LLM. & \sk{1}--\sk{3}, \sk{6} &
Test de integración con mock del LLM \\
\addlinespace
\rf{03} & Capturar audio del candidato y derivar métricas de fluidez
(palabras por minuto, palabras-muletilla, pausas, ratio
silencio/habla). & \sk{1}--\sk{3} & Test unitario sobre clip
sintético \\
\addlinespace
\rf{04} & Capturar video del candidato y derivar métricas de
contacto visual, postura y gestualidad mediante MediaPipe en el
navegador. & \sk{1}--\sk{3} & Test con grabación de prueba \\
\addlinespace
\rf{05} & Mostrar las métricas en tiempo real como ``aura luminosa''
alrededor del avatar, con latencia $\le 500\,\mathrm{ms}$. &
\sk{1}--\sk{3} & Test de rendimiento + visual regression \\
\addlinespace
\rf{06} & Ofrecer comandos de voz que reemplazan acciones de UI
(``pausa'', ``repite la pregunta'', ``terminar entrevista''). &
\sk{1}--\sk{6} & Test de NLU sobre frases de prueba \\
\addlinespace
\rf{07} & Generar al final de cada sesión un resumen y un plan de
mejora priorizado, basado en las métricas de la sesión y el historial. &
\sk{1}--\sk{2}, \sk{9} & Test de prompt $+$ snapshot del JSON
generado \\
\addlinespace
\rf{08} & Mantener un historial longitudinal de sesiones por
competencia, con tendencia visible. & \sk{1}--\sk{2} &
Test de persistencia en BD \\
\addlinespace
\rf{09} & Calibrar el nivel del candidato (junior/mid/senior) tras 3
sesiones, ajustando dificultad de futuras preguntas. & \sk{1}--\sk{2} &
Test sobre dataset de sesiones simuladas \\
\addlinespace
\rf{10} & Permitir crear y unirse a una sesión \textit{peer} en vivo vía
WebRTC, con un usuario en rol candidato y otro en rol observador. &
\sk{1}, \sk{3} & Test E2E con dos navegadores headless \\
\addlinespace
\rf{11} & Permitir al observador anclar comentarios en \textit{timestamps}
específicos de la grabación. & \sk{3}, \sk{9} & Test E2E del panel \\
\addlinespace
\rf{12} & Intercambiar roles candidato $\leftrightarrow$ observador en
cualquier momento de la sesión \textit{peer}. & \sk{1}, \sk{3} &
Test E2E \\
\addlinespace
\rf{13} & Asignar rangos por competencia (Aprendiz $\to$ Competente
$\to$ Avanzado $\to$ Experto) y \textit{badges} basados en hitos
medidos. & \sk{1}--\sk{3} & Test sobre reglas de promoción \\
\addlinespace
\rf{14} & Mantener ligas temáticas semanales por industria con
\textit{ranking}. & \sk{1}--\sk{3} & Test del motor de ranking \\
\addlinespace
\rf{15} & Mostrar \textit{tooltips} contextuales sobre cada icono del
aura explicando criterio y nivel de confianza. & \sk{1}--\sk{7} &
Test E2E de hover/foco \\
\addlinespace
\rf{16} & Proveer un asistente lateral LLM ligero que responda dudas
durante configuración y tras cada sesión. & \sk{1}--\sk{7} &
Test del prompt del asistente \\
\addlinespace
\rf{17} & Ofrecer modo \textit{audio-first} con descripción auditiva
del aura y resumen post-sesión por TTS. & \sk{4} & Test de TTS y de
descripción de eventos \\
\addlinespace
\rf{18} & Generar subtítulos en vivo con retardo $\le 300\,\mathrm{ms}$
y permitir entrada por texto en lugar de voz. & \sk{5} & Test de
latencia de subtítulos \\
\addlinespace
\rf{19} & Permitir operación 100\,\% por teclado y comandos de voz
para todas las acciones de UI. & \sk{6} & Test de cobertura
keyboard-only \\
\addlinespace
\rf{20} & Activar modo de baja estimulación con paleta neutra, sin
animaciones, aura simplificada y lenguaje claro. & \sk{7} &
Test visual + revisión de prompt en modo claro \\
\addlinespace
\rf{21} & Solicitar consentimiento informado granular (audio, video,
transcripción, métricas) antes de la primera sesión. & \sk{1}--\sk{7},
\sk{16} & Test de flujo de onboarding \\
\addlinespace
\rf{22} & Permitir al usuario eliminar grabaciones y exportar sus datos
(derecho ARCO). & \sk{1}--\sk{7}, \sk{16} & Test E2E de borrado y
exportación \\
\addlinespace
\rf{23} & Panel de administración para gestión de usuarios,
contenido de industrias y moderación del feed comunitario. &
\sk{12} & Test E2E del backoffice \\
\addlinespace
\rf{24} & Reportes agregados anonimizados de uso para institución
educativa. & \sk{10} & Test de agregación y anonimización \\

\bottomrule
\end{longtable}

\subsection{Requerimientos no funcionales}

\begin{longtable}{p{1.4cm} p{2.6cm} p{6.8cm} p{2.8cm}}
\toprule
\textbf{ID} & \textbf{Categoría} & \textbf{Descripción} & \textbf{Métrica / criterio} \\
\midrule
\endhead

\rnf{01} & Rendimiento & Latencia extremo a extremo del aura de
métricas. & $\le 500\,\mathrm{ms}$ p95 \\
\addlinespace
\rnf{02} & Rendimiento & Latencia de respuesta del LLM
entrevistador. & $\le 2\,\mathrm{s}$ p95 \\
\addlinespace
\rnf{03} & Rendimiento & La aplicación debe ejecutarse fluidamente
en una laptop estándar (8\,GB RAM, sin GPU) y en un móvil de gama
media (4\,GB RAM). & 30\,FPS sostenidos en pestaña activa \\
\addlinespace
\rnf{04} & Disponibilidad & Disponibilidad del servicio. & $\geq$~99\,\%
mensual fuera de mantenimiento \\
\addlinespace
\rnf{05} & Privacidad & No enviar video crudo a servicios externos;
solo métricas derivadas localmente por MediaPipe. & Auditoría de
tráfico saliente \\
\addlinespace
\rnf{06} & Privacidad & Cifrado en reposo de grabaciones y borrado
automático a los 30 días salvo indicación del usuario. & Test de
cron de retención \\
\addlinespace
\rnf{07} & Seguridad & Autenticación con OAuth/OIDC y \textit{tokens}
con expiración. & Test de expiración \\
\addlinespace
\rnf{08} & Accesibilidad & Cumplir WCAG 2.2 nivel AA. &
axe-core sin violaciones críticas; Lighthouse $\ge 95$ en
accesibilidad \\
\addlinespace
\rnf{09} & Internacionalización & UI y prompts del LLM en español
neutro; arquitectura preparada para añadir inglés. & Tests con
\textit{locales} es-PE y en-US \\
\addlinespace
\rnf{10} & Usabilidad & Un usuario nuevo debe poder iniciar su primera
sesión en $\le 2$ minutos desde el alta. & Test de usabilidad
cronometrado \\
\addlinespace
\rnf{11} & Mantenibilidad & Cobertura de tests automatizados $\geq$~70\,\%
en backend y $\geq$~60\,\% en frontend. & Reporte de cobertura en CI \\
\addlinespace
\rnf{12} & Observabilidad & Logs estructurados y trazas
distribuidas con correlación por sesión. & Test de presencia de
\texttt{traceId} \\
\addlinespace
\rnf{13} & Costo de operación & Costo medio de IA por sesión completa
de 30 min $\le \$0.20$ USD. & Métrica de telemetría agregada \\
\addlinespace
\rnf{14} & Ética / IA & Mostrar nivel de confianza y advertencia de
sesgo (acentos, estilos) en toda métrica derivada de modelos. &
Test del componente \textit{tooltip} \\
\addlinespace
\rnf{15} & Compatibilidad & Funcionar en Chrome, Firefox y Safari
versiones $-2$. & Tests E2E en matriz de navegadores \\

\bottomrule
\end{longtable}

\subsection{Trazabilidad stakeholders a requerimientos}
\begin{itemize}
  \item \sk{1}--\sk{2} (candidato general y con experiencia):
        \rf{01}--\rf{09}, \rf{13}--\rf{16}, \rf{21}--\rf{22}, todos los RNF.
  \item \sk{3} (peer entrevistador): \rf{10}--\rf{12}, \rf{15}--\rf{16}.
  \item \sk{4} (visual): \rf{15}, \rf{17}, \rf{21}--\rf{22}, \rnf{08}.
  \item \sk{5} (auditiva): \rf{15}, \rf{18}, \rf{21}--\rf{22}, \rnf{08}.
  \item \sk{6} (motora): \rf{06}, \rf{15}, \rf{19}, \rf{21}--\rf{22},
        \rnf{08}.
  \item \sk{7} (cognitiva/neurodivergente): \rf{15}, \rf{20},
        \rf{21}--\rf{22}, \rnf{08}.
  \item \sk{9} (coach): \rf{07}, \rf{11}.
  \item \sk{10} (universidad): \rf{24}, \rnf{04}, \rnf{13}.
  \item \sk{12} (administrador): \rf{23}.
  \item \sk{16} (regulador): \rf{21}--\rf{22}, \rnf{05}--\rnf{07}.
\end{itemize}

\section{Cumplimiento de los módulos UX/UI exigidos}
\label{sec:modulos}

El enunciado exige los \textbf{cuatro} módulos obligatorios y \textbf{tres
de los cuatro} módulos opcionales. Esta sección integra la respuesta del
proveedor Claude (prompt 04), seleccionada por su mayor precisión
mecanística (mapeo concreto de variables del aura, descripción del
\textit{pipeline} de roles del LLM, justificación estructurada del
módulo descartado) y por su trazabilidad explícita con el backlog de
RF/RNF de la sección~\ref{sec:requerimientos}.

\subsection{Módulos obligatorios}

\subsubsection{1. Interfaz gráfica no convencional}
El simulador implementa una sala virtual tridimensional renderizada con
WebGL/Three.js, en la que el entrevistador es un avatar estilizado
rodeado por un \textbf{aura luminosa procedural} que actúa como canal de
feedback continuo. El aura mapea las métricas multimodales a parámetros
visuales independientes:

\begin{itemize}
  \item \textbf{Color}: codifica la valencia global del desempeño
        (ámbar a verde).
  \item \textbf{Intensidad}: codifica el nivel de fluidez verbal.
  \item \textbf{Ritmo de pulsación}: codifica la cadencia detectada.
  \item \textbf{Asimetría espacial}: codifica el equilibrio entre canal
        verbal y no verbal.
\end{itemize}

\noindent La interacción ocurre dentro del entorno tridimensional sin
ventanas modales sobreimpuestas, lo que constituye una desviación del
paradigma de formulario web convencional. La visualización mantiene
$\ge 30$\,fps en hardware de referencia y degrada a una representación
bidimensional simplificada cuando se detecta GPU insuficiente.
\textbf{Cubre} \rf{05}, \rf{15}, \rnf{03}.

\subsubsection{2. Interfaz de voz}
Canal de voz bidireccional cuasi-real:

\begin{itemize}
  \item Las respuestas del entrevistador LLM se sintetizan mediante TTS
        (modelo de Gemini con voz seleccionable).
  \item Las intervenciones del candidato se transcriben mediante STT
        continuo con detección de fin de habla y segmentación por
        turnos.
  \item Sobre la transcripción y la señal acústica se calculan
        \textbf{en cliente} las métricas prosódicas que alimentan el
        aura (fluidez, ritmo, pausas, tono, densidad de muletillas).
  \item Una \textbf{capa de comandos de voz} permite controlar la
        sesión mediante intenciones reservadas (``pausar'',
        ``repetir pregunta'', ``siguiente'', ``terminar entrevista'')
        procesadas por un reconocedor de gramática restringida
        independiente del canal sustantivo.
\end{itemize}

\noindent La latencia mediana de TTS y de aparición de subtítulos se
mantiene por debajo del umbral fijado por \rnf{01}--\rnf{02}.
\textbf{Cubre} \rf{02}, \rf{06}, \rf{17}--\rf{19}, \rnf{01}--\rnf{02}.

\subsubsection{3. Integración con uno o varios LLM}
El sistema integra \textbf{tres roles diferenciados} sobre la API de
Gemini, cada uno con su propio prompt sistémico, contexto y memoria
acotada:

\begin{itemize}
  \item \textbf{Entrevistador}: genera preguntas adaptadas a industria,
        rol y nivel del candidato; reformula en función del historial
        conversacional de la sesión.
  \item \textbf{Coach}: opera en modo asincrónico al cierre de la
        sesión; consume la transcripción anonimizada y las métricas
        agregadas para producir un plan de mejora estructurado por
        competencias.
  \item \textbf{Moderador}: actúa exclusivamente en sesiones de peer
        mock; observa la transcripción cruzada de ambos pares y emite
        intervenciones breves cuando detecta desviaciones del
        protocolo o silencios prolongados.
\end{itemize}

\noindent Ningún rol recibe el video crudo del candidato: recibe
únicamente los descriptores numéricos derivados del \textit{pipeline}
de MediaPipe en cliente. \textbf{Cubre} \rf{02}, \rf{07}, \rf{16},
\rnf{05}, \rnf{14}.

\subsubsection{4. Ayuda contextual para el usuario}
Cuatro capas observables:

\begin{enumerate}
  \item \textbf{Tour guiado} al primer ingreso, navegable con teclado,
        que destaca cada zona funcional de la sala virtual.
  \item \textbf{Etiquetas emergentes} (\textit{tooltips}) accesibles
        por foco y por puntero sobre todos los controles, con
        descripción breve y atajo de teclado asociado.
  \item \textbf{Asistente coach-on-demand}, invocable mediante el
        comando ``ayuda'' o tecla dedicada, que responde dudas sobre
        el funcionamiento del simulador (no sobre el contenido
        evaluativo de la entrevista) consultando una base documental
        indexada.
  \item \textbf{Inferencia de necesidad de ayuda} por inactividad o
        errores repetidos: si el usuario permanece sin acción durante
        un umbral configurable o reintenta el mismo comando con falla,
        el sistema sugiere la sección de ayuda relevante.
\end{enumerate}

\noindent \textbf{Cubre} \rf{15}--\rf{16}.

\subsection{Módulos opcionales seleccionados (3 de 4)}

\subsubsection{Personalización continua}
\begin{itemize}
  \item Historial conversacional persistente, consultable como línea
        de tiempo de sesiones con métricas, transcripción y reporte
        (\rf{08}).
  \item Plan adaptativo recalculado al cierre de cada sesión: el
        módulo \textit{Coach} prioriza las competencias con menor
        desempeño en las últimas $N$ sesiones y selecciona los
        siguientes ejercicios en función de esa priorización (\rf{07}).
  \item Perfil de industria y rol declarado por el usuario y refinado
        a partir del CV cargado y de las respuestas previas; condiciona
        el banco de preguntas y el vocabulario técnico esperado por el
        entrevistador (\rf{01}).
  \item \textbf{Calibración progresiva del nivel} mediante un modelo
        de estimación de habilidad por competencia (similar a IRT) que
        sube o baja la dificultad de la siguiente pregunta según la
        calidad de la respuesta previa, con nivel inicial fijado por
        una sesión diagnóstica (\rf{09}).
  \item Persistencia de preferencias de interfaz, accesibilidad y
        configuración del aura entre sesiones y dispositivos, con
        sincronización por cuenta autenticada (\rnf{04}).
\end{itemize}

\subsubsection{Interactividad con otros usuarios}
\begin{itemize}
  \item Peer mock vía WebRTC con señalización por servidor propio y
        media P2P entre pares; sala con dos roles activos
        (candidato y observador) y un moderador LLM compartido
        (\rf{10}).
  \item Panel del observador con métricas multimodales del candidato
        en vivo (con consentimiento explícito previo del candidato
        registrado en sesión) y formulario estructurado para emitir
        feedback por competencia (\rf{11}).
  \item Comentarios anclados al \textit{timestamp} de la grabación
        local de la sesión, navegables por marcador y editables por el
        observador hasta el cierre (\rf{11}).
  \item Intercambio de roles mediante un único control que invierte
        candidato y observador y reinicia el contador de la nueva
        entrevista, preservando el historial de la sesión anterior
        como adjunto consultable (\rf{12}).
  \item \textbf{Modo panel} con hasta tres observadores simultáneos
        sobre un único candidato, con feedback agregado al cierre y
        resolución de conflictos por promedio ponderado por
        antigüedad del observador.
  \item \textbf{Feedback cruzado asincrónico}: solicitar revisión de
        una entrevista grabada por un par seleccionado del directorio,
        con plazo configurable y notificación al solicitante cuando el
        feedback queda disponible.
\end{itemize}

\subsubsection{Gamificación}
\begin{itemize}
  \item Rangos por competencia (comunicación verbal, lenguaje
        corporal, contenido técnico, manejo del estrés) calculados
        como agregaciones móviles sobre las últimas $N$ sesiones, con
        umbrales fijos y observables en el perfil (\rf{13}).
  \item \textit{Badges} otorgados por logros verificables y
        deterministas (p.\,ej. completar diez sesiones consecutivas,
        alcanzar puntaje umbral en una competencia, completar un peer
        mock como observador), con criterios documentados y reversibles
        ante anulación de la sesión correspondiente (\rf{13}).
  \item Ligas temáticas agrupadas por industria y rol, con
        clasificación semanal basada en \emph{mejora relativa} del
        usuario respecto a su propia línea base, no en comparación
        absoluta entre usuarios, para evitar penalizar a principiantes
        (\rf{14}).
  \item Retos grupales con objetivo común (acumulación colectiva de
        horas de práctica o de peer mocks completados en un periodo) y
        progreso visible para todos los miembros del grupo.
\end{itemize}

\subsection{Módulo opcional descartado: interfaces hápticas}

Se descarta la inclusión de interfaces hápticas en el alcance del MVP
por cuatro razones convergentes:

\begin{enumerate}
  \item \textbf{Dominio}: la entrevista laboral simulada no requiere
        intercambio táctil; la totalidad del valor evaluativo y
        formativo se transmite por canales auditivo, visual y textual,
        y no existe evidencia documentada de mejora de desempeño por
        háptica en este dominio.
  \item \textbf{Hardware}: la única superficie háptica realmente
        accesible en el contexto de uso esperado es la vibración del
        móvil, un canal de muy baja resolución expresiva que además no
        está disponible en escritorio (entorno principal del simulador).
  \item \textbf{Inmersión}: una vibración durante la entrevista
        interrumpiría la simulación de un contexto formal y se
        percibiría como notificación intrusiva en lugar de información
        integrada, contradiciendo la metáfora del
        \textit{Warachikuy}\footnote{El \textit{Warachikuy} (del quechua
        \textit{wara}, taparrabo otorgado en la mayoría de edad) era el
        rito de iniciación a la adultez en el Tahuantinsuyu, en el que
        los jóvenes superaban pruebas rigurosas de valor y destreza
        para recibir el reconocimiento de la comunidad y de la
        autoridad. La ceremonia, declarada de interés cultural por el
        Ministerio de Cultura del Perú, se reescenifica anualmente en
        la explanada de Saqsayhuamán. El proyecto adopta esta metáfora
        andina —en lugar de la del \textit{dojo} oriental utilizada en
        la propuesta inicial— porque captura con mayor precisión el
        carácter de tránsito profesional, examen estructurado y
        reconocimiento al cierre que define a una entrevista laboral.}.
  \item \textbf{Recursos}: el equipo prefiere profundizar los tres
        módulos opcionales seleccionados, donde existen requerimientos
        automatizables y verificables, antes que dispersar esfuerzo en
        una capa háptica cuya verificación automática sería difícil.
\end{enumerate}

\noindent La decisión queda sujeta a revisión en una fase posterior si
se incorpora un caso de uso específico (por ejemplo, alertas hápticas
para usuarios con discapacidad auditiva durante peer mock) que
justifique el costo de implementación.

\section{Especificación detallada de casos de uso}
\label{sec:casos}

Cada caso de uso sigue el formato extendido de Cockburn (actor, objetivo,
precondiciones, flujo principal, alternativos, postcondiciones) y se enlaza
con los requerimientos que satisface. La cobertura completa se resume en
la matriz de trazabilidad al final de la sección.

\subsection{\cu{01}: Realizar entrevista individual}
\begin{cajaCriterio}
\textbf{Actor primario:} Candidato (\sk{1}, \sk{2}). \\
\textbf{Actores secundarios:} LLM Entrevistador, LLM Coach. \\
\textbf{Objetivo:} Practicar una entrevista de 30 min y obtener un
resumen con plan de mejora. \\
\textbf{Precondiciones:} Usuario autenticado; consentimiento de audio,
video y métricas otorgado; navegador con permisos de cámara y micrófono. \\
\textbf{Postcondiciones (éxito):} Sesión persistida; métricas calculadas;
plan de mejora generado y mostrado al usuario.
\end{cajaCriterio}

\paragraph{Flujo principal}
\begin{enumerate}
  \item El candidato selecciona industria, rol y nivel objetivo.
  \item El sistema valida los permisos del navegador y carga la sala
        virtual.
  \item El avatar entrevistador saluda por voz y formula la primera
        pregunta.
  \item El candidato responde por voz; MediaPipe deriva métricas en el
        navegador y las envía al backend.
  \item El aura luminosa actualiza los indicadores con latencia
        $\le 500\,\mathrm{ms}$.
  \item El entrevistador formula la siguiente pregunta adaptándose a la
        respuesta previa. Se repite hasta completar la sesión o hasta
        que el candidato diga ``terminar entrevista''.
  \item El LLM Coach genera el resumen y plan de mejora; el sistema lo
        muestra al candidato y lo persiste.
\end{enumerate}

\paragraph{Flujos alternativos}
\begin{itemize}
  \item \textbf{4a.} Si MediaPipe no detecta rostro durante $>10\,\mathrm{s}$,
        el sistema pausa y pide reencuadre (sin penalizar postura).
  \item \textbf{6a.} Si el candidato dice ``pausa'', el sistema detiene
        captura y temporizador hasta el comando ``continuar''.
  \item \textbf{6b.} Si el candidato dice ``repite la pregunta'', el
        avatar la reformula sin contarlo como reintento.
  \item \textbf{2a / 7a.} Si falla la API del LLM, el sistema cae en una
        cola de reintentos y avisa al usuario sin perder la sesión.
\end{itemize}

\paragraph{Cubre} \rf{01}--\rf{07}, \rf{15}, \rnf{01}--\rnf{02},
\rnf{05}, \rnf{14}.

\subsection{\cu{02}: Configurar perfil y preferencias de accesibilidad}
\begin{cajaCriterio}
\textbf{Actor primario:} Cualquier usuario primario (\sk{1}--\sk{7}). \\
\textbf{Objetivo:} Establecer industria preferida, idioma, paciencia del
entrevistador, modo de baja estimulación, modo audio-first, contraste,
tipografía, declaración de condiciones que reemplazan métricas. \\
\textbf{Precondiciones:} Usuario autenticado.
\end{cajaCriterio}

\paragraph{Flujo principal}
\begin{enumerate}
  \item El usuario abre ``Mi perfil''.
  \item Selecciona o ajusta cada preferencia; la UI muestra una vista
        previa en vivo del impacto.
  \item Guarda; las preferencias se aplican inmediatamente a sesiones
        futuras y al asistente lateral.
\end{enumerate}

\paragraph{Cubre} \rf{17}--\rf{20}, \rnf{08}--\rnf{09}.

\subsection{\cu{03}: Realizar mock interview entre pares (peer)}
\begin{cajaCriterio}
\textbf{Actores primarios:} Candidato (\sk{1}) y Peer Observador (\sk{3}). \\
\textbf{Objetivo:} Practicar una entrevista con feedback humano sin
salir de la aplicación. \\
\textbf{Precondiciones:} Ambos usuarios autenticados; uno de ellos creó
una sala con código compartible; ambos consintieron grabación.
\end{cajaCriterio}

\paragraph{Flujo principal}
\begin{enumerate}
  \item El candidato crea sala; el sistema emite un código.
  \item El peer entra con el código; WebRTC establece la conexión.
  \item El candidato responde a preguntas que pueden venir del LLM
        entrevistador, del peer, o de ambos según el modo.
  \item El peer ancla comentarios en la línea de tiempo durante o
        después de cada respuesta.
  \item Al terminar, los usuarios pueden intercambiar roles.
  \item El sistema genera dos resúmenes (uno por rol) y los persiste.
\end{enumerate}

\paragraph{Flujos alternativos}
\begin{itemize}
  \item \textbf{2a.} Si la conexión WebRTC falla, el sistema reintenta
        con servidor TURN; si persiste, ofrece modo asincrónico
        (grabar y compartir).
\end{itemize}

\paragraph{Cubre} \rf{10}--\rf{12}, \rf{21}, \rnf{04}, \rnf{15}.

\subsection{\cu{04}: Consultar historial y plan de mejora}
\begin{cajaCriterio}
\textbf{Actor primario:} Candidato (\sk{1}, \sk{2}). \\
\textbf{Actor secundario:} Coach o mentor (\sk{9}) si el usuario lo
autorizó. \\
\textbf{Objetivo:} Revisar la evolución por competencia, las sesiones
pasadas y el plan de mejora vigente.
\end{cajaCriterio}

\paragraph{Flujo principal}
\begin{enumerate}
  \item El usuario abre ``Mi progreso''.
  \item Visualiza tendencias por competencia (gráfico) y \textit{badges}
        obtenidos.
  \item Selecciona una sesión específica para revisar la grabación,
        las métricas y los comentarios.
  \item Pide al asistente lateral que profundice en un punto del plan
        (``dame ejemplos para reducir muletillas''); el LLM responde con
        contexto del historial.
\end{enumerate}

\paragraph{Cubre} \rf{07}--\rf{09}, \rf{16}, \rnf{14}.

\subsection{\cu{05}: Gestionar gamificación (rangos, badges, ligas)}
\begin{cajaCriterio}
\textbf{Actor primario:} Candidato (\sk{1}--\sk{3}). \\
\textbf{Objetivo:} Inscribirse en una liga, ver ranking, recibir
\textit{badges} al cumplir hitos.
\end{cajaCriterio}

\paragraph{Flujo principal}
\begin{enumerate}
  \item El usuario explora ligas activas filtrando por industria.
  \item Se inscribe en la liga correspondiente.
  \item Cada sesión válida actualiza su puntaje semanal.
  \item Al cumplir el hito de un \textit{badge} (ej. 3 sesiones con
        cero muletillas en 5 min), el sistema lo otorga y notifica.
  \item Al cierre semanal, se publica el ranking y el usuario puede
        ver su posición.
\end{enumerate}

\paragraph{Cubre} \rf{13}--\rf{14}, \rnf{04}.

\subsection{\cu{06}: Gestionar consentimiento, datos y privacidad}
\begin{cajaCriterio}
\textbf{Actor primario:} Cualquier usuario primario (\sk{1}--\sk{7}). \\
\textbf{Actor regulador:} Autoridad de protección de datos (\sk{16}). \\
\textbf{Objetivo:} Otorgar/revocar consentimientos, eliminar grabaciones,
exportar datos.
\end{cajaCriterio}

\paragraph{Flujo principal}
\begin{enumerate}
  \item En el primer ingreso, el sistema solicita consentimiento granular
        (audio, video, transcripción, métricas, uso comunitario).
  \item En cualquier momento, el usuario puede revocar un consentimiento
        específico desde ``Mi privacidad''.
  \item El usuario solicita borrado de una grabación; el sistema lo
        ejecuta y emite confirmación.
  \item El usuario solicita exportación de datos; el sistema genera un
        archivo cifrado y le envía un enlace temporal.
\end{enumerate}

\paragraph{Cubre} \rf{21}--\rf{22}, \rnf{05}--\rnf{07}.

\subsection{\cu{07}: Administrar la plataforma}
\begin{cajaCriterio}
\textbf{Actor primario:} Administrador del sistema (\sk{12}). \\
\textbf{Objetivo:} Gestionar usuarios, contenido de industrias, ligas y
moderación del feed comunitario.
\end{cajaCriterio}

\paragraph{Flujo principal}
\begin{enumerate}
  \item El administrador entra al \textit{backoffice}.
  \item Crea, edita o desactiva ligas, industrias o usuarios.
  \item Revisa el feed comunitario; modera o elimina contenido reportado.
\end{enumerate}

\paragraph{Cubre} \rf{23}.

\subsection{\cu{08}: Generar reportes institucionales}
\begin{cajaCriterio}
\textbf{Actor primario:} Universidad / oficina de empleabilidad (\sk{10}). \\
\textbf{Objetivo:} Obtener métricas agregadas anonimizadas para reportes
de empleabilidad.
\end{cajaCriterio}

\paragraph{Flujo principal}
\begin{enumerate}
  \item Un usuario institucional autenticado accede al panel de reportes.
  \item Filtra por industria, rango de fechas y rangos por competencia.
  \item Descarga un reporte en CSV/PDF con datos agregados (nunca
        identificables).
\end{enumerate}

\paragraph{Cubre} \rf{24}, \rnf{05}.

\subsection{Matriz de trazabilidad caso de uso a requerimiento}

\begin{longtable}{p{1.6cm} p{12cm}}
\toprule
\textbf{CU} & \textbf{Requerimientos cubiertos} \\
\midrule
\cu{01} & \rf{01}--\rf{07}, \rf{15}; \rnf{01}, \rnf{02}, \rnf{05}, \rnf{14} \\
\cu{02} & \rf{17}--\rf{20}; \rnf{08}, \rnf{09} \\
\cu{03} & \rf{10}--\rf{12}, \rf{21}; \rnf{04}, \rnf{15} \\
\cu{04} & \rf{07}--\rf{09}, \rf{16}; \rnf{14} \\
\cu{05} & \rf{13}--\rf{14}; \rnf{04} \\
\cu{06} & \rf{21}--\rf{22}; \rnf{05}--\rnf{07} \\
\cu{07} & \rf{23} \\
\cu{08} & \rf{24}; \rnf{05} \\
\bottomrule
\end{longtable}

\noindent Los RNF transversales \rnf{03}, \rnf{08}, \rnf{10}--\rnf{13} se
verifican a nivel de toda la aplicación (rendimiento global, accesibilidad
en todas las vistas, mantenibilidad, observabilidad y costo) y no se
asignan a un caso de uso individual.

\section{Proveedores de IA generativa y prompts utilizados}
\label{sec:prompts}

La tarea LEC06 exige el uso de \textbf{al menos dos proveedores de IA
generativa}, indicar los \textit{prompts} optimizados y justificar por
qué los resultados elegidos son los mejores. El equipo empleó los dos
siguientes, ambos en sus interfaces web:

\begin{itemize}
  \item \textbf{Google Gemini} (versión web, modelo Gemini 2.5).
        Priorizado en tareas que requerían \emph{volumen y cobertura}:
        listado amplio de \textit{stakeholders} y guía operativa de
        entrevistas con representantes de usuarios.
  \item \textbf{Anthropic Claude} (versión web, familia Claude 4.X).
        Priorizado en tareas que requerían \emph{rigor estructural,
        consistencia de formato y sensibilidad ética}: funcionalidad
        para usuarios con discapacidad, cumplimiento de módulos UX/UI
        y casos de uso en formato Cockburn extendido.
\end{itemize}

Las respuestas crudas se versionaron en
\texttt{prompts/respuestas/} para auditoría: cada salida puede
contrastarse contra el prompt que la originó y contra la versión
integrada en el informe.

\subsection{Estrategia general}

Toda conversación con cada IA comenzó con un \textbf{contexto base}
idéntico (\texttt{prompts/00\_contexto\_base.md}) que fija dominio,
restricciones de estilo, idioma, lenguaje centrado en la persona y el
principio de \emph{automatización obligatoria}: cualquier
funcionalidad o requerimiento propuesto debe ser verificable por test,
telemetría o validación de UI.

Para cada uno de los puntos exigidos por la tarea se eligió el
proveedor con mayor probabilidad de producir el mejor resultado, y los
resultados se compilaron manualmente en este informe sin copiar
respuestas de IA en bruto. Las divergencias entre la respuesta del
proveedor y el contenido final del informe se documentan en cada
subsección.

\subsection{Prompt 1 --- Stakeholders (Gemini)}
\paragraph{Archivos} prompt: \texttt{prompts/01\_stakeholders\_gemini.md};
respuesta: \texttt{prompts/respuestas/gemini\_01\_stakeholders.md}.

\paragraph{Por qué Gemini} Su tendencia a producir listas amplias y a
incorporar actores institucionales y regulatorios sin que se le
recuerde explícitamente facilita maximizar cobertura desde el primer
borrador.

\paragraph{Criterio de selección y resultado} La respuesta de Gemini
cubrió los cuatro perfiles con discapacidad como \textit{stakeholders}
primarios separados, identificó la Ley
N\textsuperscript{\textordmasculine}~29733 como marco regulatorio peruano
relevante, y nombró concretamente la Oficina de Empleabilidad UNI y
los proveedores cloud (Render, Fly.io, Vercel). Estos descriptores
específicos se integraron en la sección~\ref{sec:stakeholders}.

\paragraph{Divergencia documentada} Gemini propuso 14
\textit{stakeholders} y omitió cuatro roles que el equipo consideró
relevantes: \textit{peer} entrevistador como rol activo, reclutador
como consultor del proyecto, coach/mentor y administrador del sistema
diferenciado del equipo de desarrollo. La sección 01 los agrega
explícitamente y mantiene su numeración \sk{1}--\sk{17} como
autoritativa, con tabla de equivalencia hacia la numeración de Gemini.

\paragraph{Incidencia detectada} En la ejecución del 29-04-2026, la
salida del prompt 02 fue idéntica a la del prompt 01 en la sesión web
de Gemini (probable reutilización de contexto en la pestaña). El
equipo decidió no reejecutar el prompt 01 y delegar el contenido
sustantivo del punto 2 (funcionalidad para discapacidad) a Claude, que
es donde realmente se requería el rigor de WCAG.

\subsection{Prompt 2 --- Funcionalidad para discapacidad (Claude)}
\paragraph{Archivos} prompt: \texttt{prompts/02\_discapacidad\_claude.md};
respuesta: \texttt{prompts/respuestas/claude\_02\_discapacidad.md}.

\paragraph{Por qué Claude} La sección requiere lenguaje centrado en la
persona, citas correctas a WCAG 2.2 nivel AA y atención a casos donde
una métrica multimodal puede penalizar injustamente al usuario.

\paragraph{Criterio de selección y resultado} La respuesta de Claude
cumplió los cuatro criterios definidos a priori: (i) propuso métricas
sustitutas para cada caso de penalización injusta (postura
$\to$ presencia, contacto visual $\to$ orientación de voz hacia el
micrófono, fluidez verbal $\to$ coherencia y densidad léxica, postura
$\to$ estabilidad del encuadre facial); (ii) citó criterios WCAG
concretos por número (2.1.1, 1.4.3, 1.4.4, 1.2.4, 1.4.13, 1.4.12,
2.5.8, 2.2.1, 2.1.4, 2.5.7, 2.5.4, 2.3.3, 3.2.5, 4.1.2); (iii) usó
lenguaje centrado en la persona en todo el texto; (iv) cerró con un
párrafo de verificación que enumera herramientas (\texttt{axe-core},
Lighthouse, Playwright, panel de telemetría) y cuatro hitos por perfil.
La respuesta se integró íntegramente en la sección~\ref{sec:discapacidad}
con ajustes mínimos de formato LaTeX.

\subsection{Prompt 3 --- Requerimientos funcionales y no funcionales}
\paragraph{Estado} Pendiente de ejecutar formalmente con un proveedor.
El equipo elaboró la sección~\ref{sec:requerimientos} a partir del
\textit{backlog} consensuado en las reuniones de diseño, validado
contra los \textit{stakeholders} de la sección~\ref{sec:stakeholders}
y contra los módulos UX/UI de la sección~\ref{sec:modulos}. El prompt
preparado para Gemini (\texttt{prompts/03\_requerimientos\_gemini.md})
está listo para ejecución en una iteración posterior si el profesor
lo requiere; los criterios de selección ya están definidos en el
prompt.

\subsection{Prompt 4 --- Cumplimiento de módulos UX/UI (Claude)}
\paragraph{Archivos} prompt: \texttt{prompts/04\_modulos\_claude.md};
respuesta: \texttt{prompts/respuestas/claude\_04\_modulos.md}.

\paragraph{Por qué Claude} La sección requiere coherencia
argumentativa entre lo que cada módulo es, cómo lo cumple el
simulador y qué requerimientos satisface.

\paragraph{Criterio de selección y resultado} Se seleccionó la
respuesta de Claude por: (i) precisión mecanística en el mapeo del
aura (color $=$ valencia, intensidad $=$ fluidez, ritmo de pulsación
$=$ cadencia, asimetría espacial $=$ equilibrio verbal/no verbal);
(ii) descripción del \textit{pipeline} de tres roles del LLM con
contexto y memoria acotada propios; (iii) justificación de la
exclusión de hápticas con cuatro argumentos diferenciados (dominio,
hardware, inmersión, recursos), no agrupados en un único párrafo
difuso; (iv) cuatro capas observables de ayuda contextual, no la
clásica enumeración tooltip/asistente. Estos elementos se integraron
en la sección~\ref{sec:modulos}.

\paragraph{Divergencia documentada} Claude usó una numeración de
RF/RNF distinta de la canónica del informe (RF-21..RF-60). Las
referencias en la sección~\ref{sec:modulos} fueron remapeadas a la
numeración \rf{01}--\rf{24} de la sección~\ref{sec:requerimientos}.

\subsection{Prompt 5 --- Casos de uso (Claude)}
\paragraph{Archivos} prompt: \texttt{prompts/05\_casos\_uso\_claude.md};
respuesta: \texttt{prompts/respuestas/claude\_05\_casos\_uso.md}.

\paragraph{Por qué Claude} Para ocho casos de uso largos, la
consistencia entre todos importa más que la cobertura adicional.
Claude respeta plantillas extensas sin fugarse del esquema y, cuando
detecta que falta un RF, lo señala como nota separada al final en
lugar de inventarlo en medio.

\paragraph{Criterio de selección y resultado} Se seleccionó la
respuesta de Claude porque: (i) produjo exactamente los ocho casos de
uso pedidos con la numeración solicitada; (ii) cada CU incluyó al
menos un flujo alternativo documentado; (iii) cerró con dos notas
explícitas sobre cobertura: tres RNF transversales que no encajan en
un CU específico (compatibilidad de navegador, cobertura de pruebas,
consumo de memoria) y dos posibles RF faltantes (matchmaking previo a
peer mock, persistencia de preferencias de accesibilidad entre
dispositivos). Estas observaciones se incorporaron como notas finales
en la sección~\ref{sec:casos}. El cuerpo de los CU del informe usa la
numeración \rf{}/\rnf{} canónica.

\paragraph{Divergencia documentada} Como en el prompt 4, la respuesta
de Claude empleó una numeración propia de RF/RNF (RF-46, RF-60, etc.)
que sugería un \textit{backlog} ampliado. El informe mantiene la
numeración \rf{01}--\rf{24} de la sección~\ref{sec:requerimientos}
como autoritativa; las observaciones de Claude sobre RF faltantes
quedan como insumo para una iteración futura del \textit{backlog}.

\subsection{Prompt 6 --- Guía de entrevistas con usuarios finales (Gemini)}
\paragraph{Archivos} prompt:
\texttt{prompts/06\_entrevistas\_gemini.md}; respuesta:
\texttt{prompts/respuestas/gemini\_06\_entrevistas.md}.

\paragraph{Por qué Gemini} La generación de preguntas variadas para
siete perfiles distintos en una sola pasada se beneficia de la
amplitud de Gemini.

\paragraph{Criterio de selección y resultado} La respuesta cumplió los
criterios: (i) preguntas abiertas en presente sin sesgar la respuesta;
(ii) tiempo total acotado a 45--60 minutos; (iii) bloques específicos
por perfil con discapacidad, no genéricos; (iv) plantilla de
consentimiento informado lista para imprimir. Se integró
íntegramente en la sección~\ref{sec:entrevistas} como guía operativa
para que el equipo ejecute las entrevistas reales.

\subsection{Resumen del flujo de trabajo}

\begin{enumerate}
  \item Pegar el contexto base en una conversación nueva en cada IA.
  \item Esperar la confirmación ``Listo'' antes de enviar el prompt
        específico.
  \item Pegar el prompt del punto correspondiente.
  \item Versionar la respuesta cruda en
        \texttt{prompts/respuestas/<proveedor>\_<prompt>\_<topico>.md}.
  \item Aplicar los criterios de selección descritos en cada
        subsección al resultado.
  \item Si el resultado falla algún criterio, repetir con prompt
        ajustado y conservar las dos versiones.
  \item Compilar manualmente los mejores fragmentos en este informe
        LaTeX y documentar las divergencias.
\end{enumerate}

\subsection*{Apéndice A --- Restricciones del contexto base}
\begin{enumerate}
  \item Español neutro académico, sin emojis ni coloquialismos.
  \item Listas exhaustivas y mutuamente excluyentes.
  \item No inventar datos numéricos; marcar como
        \texttt{(referencial)}.
  \item Toda funcionalidad propuesta debe ser automatizable.
  \item Lenguaje centrado en la persona y referencia a WCAG 2.2 AA.
  \item Salida lista para pegar en LaTeX.
\end{enumerate}

\section{Entrevistas con representantes de usuarios finales}
\label{sec:entrevistas}

La nota final del enunciado establece textualmente:
\textit{``Se espera que en la definición formal de su proyecto (PC02)
se tomen en cuenta estos y otros resultados, incluyendo resultado de
entrevistas con representantes de sus usuarios finales''}.

Esta sección presenta la \textbf{guía operativa} elaborada por el
equipo a partir del prompt 06 (Gemini, versionado en
\texttt{prompts/respuestas/gemini\_06\_entrevistas.md}) y deja
explícitas las acciones que el equipo ejecutará para completar el
trabajo con datos reales.

\begin{cajaPrompt}[title={Estado actual --- pendiente de ejecución de campo}]
La guía operativa (perfiles, criterios de selección, bloques de
preguntas, métricas a registrar y plantilla de consentimiento) está
lista. Falta el trabajo de campo: contactar entrevistados, ejecutar
las sesiones, transcribir y consolidar hallazgos. La consolidación
final se incorporará en este mismo archivo en una iteración posterior
del informe.
\end{cajaPrompt}

\subsection{Perfil 1 --- Estudiante de últimos ciclos (\sk{1})}

\paragraph{Objetivo} Identificar puntos de dolor en la transición
académica-laboral y validar la aceptación de una interfaz de
entrenamiento basada en IA.

\paragraph{Criterios de selección}
\begin{itemize}
  \item Estudiante matriculado en 8\textsuperscript{vo},
        9\textsuperscript{no} o 10\textsuperscript{mo} ciclo de una
        carrera de ingeniería.
  \item Sin experiencia previa en prácticas preprofesionales o con
        máximo una experiencia técnica.
  \item Búsqueda activa de empleo declarada para los próximos 6 meses.
  \item Usuario recurrente de herramientas de videoconferencia
        (Zoom, Meet, Teams).
\end{itemize}

\paragraph{Bloques de preguntas}
\begin{itemize}
  \item \textbf{Contexto y experiencia previa}:
        \begin{itemize}
          \item ¿Cómo ha sido su proceso de preparación para las
                primeras postulaciones laborales?
          \item ¿Qué aspectos de una entrevista presencial o virtual le
                generan mayor incertidumbre?
        \end{itemize}
  \item \textbf{Expectativas frente a la IA}:
        \begin{itemize}
          \item ¿Qué valor esperaría obtener de un evaluador que no
                es humano?
          \item ¿Cómo reaccionaría si el nivel de dificultad de las
                preguntas cambia según sus respuestas anteriores?
        \end{itemize}
  \item \textbf{Metáfora del dojo y el aura}:
        \begin{itemize}
          \item Al imaginar un entorno de ``entrenamiento constante''
                tipo dojo, ¿qué elementos visuales consideraría
                motivadores?
          \item ¿De qué manera interpretaría un halo de luz que cambia
                de color mientras usted habla?
        \end{itemize}
  \item \textbf{Privacidad y ética}:
        \begin{itemize}
          \item ¿Qué condiciones exigiría para permitir que un
                algoritmo analice sus gestos faciales?
          \item ¿Cómo preferiría que el sistema le informe sobre el
                uso de sus datos tras finalizar la sesión?
        \end{itemize}
\end{itemize}

\paragraph{Métricas a registrar}
\begin{itemize}
  \item Frecuencia de mención de términos relacionados con
        ``ansiedad'' o ``nerviosismo''.
  \item Tiempo de respuesta ante la propuesta del aura
        (latencia en la reacción).
  \item Recuento de adjetivos positivos/negativos asociados al
        feedback automatizado.
\end{itemize}

\subsection{Perfil 2 --- Recién egresado (\sk{2})}

\paragraph{Objetivo} Evaluar la utilidad de simulaciones de alta
complejidad y la efectividad del plan de mejora personalizado para la
inserción laboral real.

\paragraph{Criterios de selección}
\begin{itemize}
  \item Egresado universitario con menos de 12 meses de haber culminado
        estudios.
  \item Postulante activo con al menos 3 entrevistas realizadas en el
        último trimestre.
  \item Acceso a equipo de cómputo propio con webcam y micrófono
        funcional.
\end{itemize}

\paragraph{Bloques de preguntas}
\begin{itemize}
  \item Experiencia en procesos reales: motivos de rechazo frecuentes,
        retroalimentación recibida.
  \item Funcionalidad del plan de mejora: información indispensable,
        uso cotidiano del plan.
  \item Metáfora y gamificación: rangos/ligas, sala virtual.
  \item Privacidad y datos: almacenamiento por la universidad,
        confianza en análisis de lenguaje corporal por IA.
\end{itemize}

\paragraph{Métricas a registrar}
\begin{itemize}
  \item Sugerencias específicas para el plan de mejora.
  \item Grado de interés en gamificación (escala 1--5 observada).
  \item Palabras clave recurrentes (``empleabilidad'', ``competencia'').
\end{itemize}

\subsection{Perfil 3 --- Coach de carrera o reclutador (\sk{8}, \sk{9})}

\paragraph{Objetivo} Contrastar las métricas automáticas del sistema
con los estándares profesionales de evaluación de talento.

\paragraph{Criterios de selección}
\begin{itemize}
  \item Profesional con más de 3 años de experiencia en reclutamiento
        IT o coaching laboral.
  \item Experiencia previa evaluando perfiles \textit{junior} o
        practicantes de ingeniería.
  \item Conocimiento básico de IA aplicada a RRHH.
\end{itemize}

\paragraph{Bloques de preguntas}
\begin{itemize}
  \item Criterios de evaluación: indicadores no verbales,
        progresión de dificultad.
  \item Validación de la herramienta: fiabilidad del análisis de
        muletillas y ritmo, sesgos.
  \item Interacción humano-IA: reacción ante respuestas incompletas,
        valor del peer mock.
  \item Privacidad institucional: garantías de seguridad, reportes
        agregados.
\end{itemize}

\paragraph{Métricas a registrar}
\begin{itemize}
  \item Términos técnicos de reclutamiento omitidos por la propuesta.
  \item Gestos de escepticismo o validación ante las métricas de
        lenguaje corporal.
  \item Tiempo dedicado a discutir la ética de la IA.
\end{itemize}

\subsection{Perfiles 4--7 --- Usuarios con discapacidad (\sk{4}--\sk{7})}

\paragraph{Objetivo} Validar la accesibilidad de la interfaz multimodal
y la eficacia de los canales alternativos de feedback.

\paragraph{Criterios de selección}
\begin{itemize}
  \item \sk{4} (visual): usuario dependiente de lector de pantalla o
        con baja visión severa.
  \item \sk{5} (auditiva): usuario con pérdida auditiva que utilice
        subtitulado o lengua de señas.
  \item \sk{6} (motora): usuario con movilidad reducida en manos que
        utilice software de asistencia.
  \item \sk{7} (cognitiva/neurodivergente): usuario diagnosticado
        dentro del espectro autista o con TDAH.
\end{itemize}

\paragraph{Bloques de preguntas (comunes y específicos)}
\begin{itemize}
  \item \textbf{Barreras de accesibilidad}: obstáculos en plataformas
        de entrevista por video actuales; interacción típica con
        asistentes de voz o comandos verbales.
  \item \textbf{Interfaz no convencional} (específico \sk{4}/\sk{5}):
        \begin{itemize}
          \item (\sk{4}) ¿Cómo esperaría recibir la información del
                aura si esta es puramente visual?
          \item (\sk{5}) ¿Qué tan útil le resulta la descripción
                textual del tono de voz del entrevistador?
        \end{itemize}
  \item \textbf{Carga cognitiva y ayuda} (específico \sk{6}/\sk{7}):
        \begin{itemize}
          \item (\sk{6}) ¿Qué dificultades prevé al interactuar con
                una interfaz que requiere respuestas en tiempo real?
          \item (\sk{7}) ¿De qué manera el avatar del entrevistador
                podría reducir la ansiedad social en lugar de
                aumentarla?
        \end{itemize}
  \item \textbf{Manejo de sensores}:
        \begin{itemize}
          \item ¿Qué inquietudes le genera que el sistema analice sus
                patrones de movimiento o voz específicos?
          \item ¿Cómo prefiere ser notificado si el sistema no logra
                procesar correctamente su imagen o audio?
        \end{itemize}
\end{itemize}

\paragraph{Métricas a registrar}
\begin{itemize}
  \item Identificación de fallos potenciales en el flujo de navegación
        mediante voz o teclado.
  \item Menciones explícitas al estándar WCAG o necesidades de
        personalización.
  \item Signos de fatiga cognitiva durante la descripción de
        funcionalidades complejas.
\end{itemize}

\subsection{Plantilla de consentimiento informado}

\begin{cajaCriterio}
\textbf{Título del proyecto}: Simulador de Entrevistas Laborales
Adaptativo (CC451 --- UNI). \\
\textbf{Investigadores}: Davila Santos, Serrano Arostegui, Poma
Navarro.

\paragraph{Propósito} Se le solicita participar en una entrevista para
recoger requisitos y percepciones sobre una plataforma de simulación
de entrevistas con IA. La información ayudará a mejorar la
accesibilidad y eficacia del sistema.

\paragraph{Procedimiento} La sesión durará entre 45 y 60 minutos. Se
registrará audio y se tomarán notas de la transcripción para análisis
académico exclusivo del curso CC451.

\paragraph{Confidencialidad} Sus datos personales serán anonimizados
mediante códigos. Las grabaciones se eliminarán al finalizar el ciclo
académico 2026-I. No se compartirá video crudo ni audio con entidades
externas al equipo de desarrollo y la cátedra.

\paragraph{Derechos} Su participación es voluntaria. Puede negarse a
responder cualquier pregunta o retirarse del proceso en cualquier
momento sin penalización alguna.

\paragraph{Consentimiento}
Yo, \rule{6cm}{0.4pt}, identificado con DNI/Código \rule{2.5cm}{0.4pt},
declaro haber sido informado sobre los objetivos del estudio y acepto
voluntariamente participar, autorizando la grabación de esta sesión
únicamente para fines de transcripción y análisis académico. \\[0.5cm]
Firma: \rule{5cm}{0.4pt} \quad Fecha: \rule{0.6cm}{0.4pt}/\rule{0.6cm}{0.4pt}/2026
\end{cajaCriterio}

\subsection{Plan de ejecución y consolidación}

\paragraph{Acciones del equipo}
\begin{enumerate}
  \item Reclutar al menos un entrevistado por cada perfil
        (\sk{1}, \sk{2}, \sk{8} o \sk{9}, y un representante por cada
        \sk{4}--\sk{7}).
  \item Realizar las entrevistas (45--60 min cada una) con
        consentimiento informado firmado.
  \item Transcribir las grabaciones y registrar hallazgos por
        entrevista.
  \item Consolidar patrones recurrentes, discrepancias y ajustes al
        \textit{backlog} de RF/RNF o casos de uso, con su justificación
        basada en los hallazgos.
  \item Incorporar la consolidación en este mismo archivo en la
        siguiente iteración del informe, con la siguiente estructura
        por entrevista: perfil del entrevistado, fecha y duración,
        modalidad, hallazgos clave por bloque temático, citas
        literales relevantes (con consentimiento explícito) y cambios
        al diseño del proyecto derivados, con referencia al RF/RNF o
        caso de uso afectado.
\end{enumerate}

\paragraph{Trazabilidad entrevistas a requerimientos} Por cada
hallazgo accionable se registrará: ID de la entrevista (E-NN),
hallazgo en una oración, requerimiento o caso de uso impactado
(\rf{N}, \rnf{N}, \cu{N}), tipo de impacto
(confirma / agrega / modifica / descarta).

\section{Sistemas y aplicaciones similares}
\label{sec:similares}

Antes de fijar el alcance del proyecto se relevaron las herramientas
disponibles en el mercado que abordan, total o parcialmente, el problema
de la preparación para entrevistas laborales. El relevamiento se centró
en cinco productos representativos del estado del arte, seleccionados
por cobertura de prensa y tracción reportada en el último año, y
distribuidos entre las dos categorías observables: simuladores
asistidos por IA (Interview Warmup, Big Interview, Yoodli, Final Round
AI) y plataformas de práctica entre pares (Pramp). Se descartaron del
relevamiento las plataformas de evaluación corporativa
(p.\,ej.\ HireVue) por estar orientadas a reclutadores y no al
candidato.

\subsection{Productos relevados}

\paragraph{Interview Warmup (Google).} Producto gratuito orientado a
candidatos de ingreso a empresas tecnológicas. Presenta preguntas en
texto y permite respuesta hablada por el navegador. Genera una
visualización de patrones lingüísticos (palabras frecuentes,
referencias profesionales, muletillas) sobre la transcripción. No
analiza video ni emite un puntaje cuantitativo y no contempla
adaptación dinámica de la dificultad ni interacción con otros usuarios.

\paragraph{Big Interview.} Suite comercial con planes individuales e
institucionales. Combina contenidos formativos extensos (videos, guías,
plantillas STAR) con un módulo de práctica donde el usuario se graba a
sí mismo respondiendo preguntas y recibe un panel posterior con
métricas básicas (\textit{eye contact}, \textit{filler words}, ritmo).
La retroalimentación se entrega como \textit{dashboard} lateral
clásico y no en tiempo real durante la entrevista. La adaptación al
candidato es limitada y depende de la curaduría manual de bancos de
preguntas por industria.

\paragraph{Yoodli.} Aplicación enfocada en la fluidez del discurso
hablado más que en el contenido técnico de la entrevista. Soporta
sesiones grabadas y modo en vivo durante reuniones de
videoconferencia, con métricas prosódicas (palabras por minuto, pausas,
muletillas, riqueza léxica) entregadas al cierre. Su análisis no
verbal se limita a la voz y no incorpora análisis postural ni de
contacto visual; tampoco habilita sesiones colaborativas entre pares.

\paragraph{Final Round AI.} Producto reciente posicionado como
\textit{interview copilot}: durante una entrevista real con un
reclutador humano, sugiere al candidato respuestas o estructuras
argumentales en una ventana lateral. El modo \textit{mock} es
secundario y replica el patrón pregunta-respuesta con feedback
posterior. El énfasis del producto está en la asistencia en vivo y no
en el entrenamiento longitudinal con plan de mejora; el componente
ético del uso del \textit{copilot} en entrevistas reales está aún en
discusión pública.

\paragraph{Pramp.} Plataforma colaborativa donde dos candidatos se
emparejan por compatibilidad de rol y se entrevistan mutuamente, sin
intervención de IA. Cada participante alterna roles candidato y
entrevistador, y se evalúan con una rúbrica predefinida. Su valor
reside en la práctica humana realista, pero la calidad del feedback
queda íntegramente en manos del par y carece de métricas objetivas
multimodales.

\subsection{Comparación funcional}

La tabla \ref{tab:similares} resume el cruce de capacidades observadas
con los ejes que estructuran este proyecto: análisis multimodal,
\textit{feedback} en tiempo real, adaptación al candidato, modalidad
colaborativa entre pares, gamificación y atención explícita a la
accesibilidad.

\begin{table}[ht]
\centering
\renewcommand{\arraystretch}{1.25}
\footnotesize
\begin{tabularx}{\linewidth}{@{}l Y Y Y Y Y c@{}}
\toprule
\textbf{Producto} &
\textbf{Análisis voz} &
\textbf{Análisis video} &
\textbf{Feedback en vivo} &
\textbf{Adaptación al nivel} &
\textbf{Práctica entre pares} &
\textbf{WCAG AA} \\
\midrule
Interview Warmup & Limitado (texto) & No & No & Limitada & No & Parcial \\
Big Interview    & Sí (post)        & Sí (post) & No & Limitada & No & Parcial \\
Yoodli           & Sí (vivo+post)   & No & Parcial & No & No & Parcial \\
Final Round AI   & Sí (vivo)        & No & Sí (sugerencias) & Limitada & No & No documentado \\
Pramp            & No               & No (peer humano) & N/A & N/A & Sí (núcleo) & No documentado \\
\midrule
\textbf{Este proyecto} &
\textbf{Sí (vivo)} &
\textbf{Sí (vivo)} &
\textbf{Sí (aura)} &
\textbf{Sí (IRT)} &
\textbf{Sí (peer mock)} &
\textbf{Sí (AA explícito)} \\
\bottomrule
\end{tabularx}
\caption{Comparación funcional con cinco productos representativos.
La fila final corresponde al simulador propuesto.}
\label{tab:similares}
\end{table}

\subsection{Posicionamiento y diferenciación}

El relevamiento permite ubicar el proyecto en una intersección poco
cubierta del espacio de soluciones existentes. La oferta actual se
distribuye en tres polos: (i) simuladores con \textit{feedback}
posterior y análisis incompleto del lenguaje no verbal (Big Interview,
Yoodli), (ii) asistentes en vivo orientados a entrevistas reales con
componente ético cuestionado (Final Round AI), y (iii) plataformas de
práctica entre pares sin métricas objetivas (Pramp). Ninguna de las
herramientas relevadas combina simultáneamente: análisis multimodal
sincrónico de voz y video con cómputo en cliente, retroalimentación
visual integrada al avatar (no como \textit{dashboard} lateral),
adaptación de dificultad por modelo de habilidad por competencia,
modalidad colaborativa entre pares con observador instrumentado, y
conformidad explícita con WCAG 2.2 nivel AA con rutas alternativas
para usuarios con discapacidad.

La propuesta del presente proyecto se diferencia en cuatro ejes
verificables. Primero, el \textit{aura luminosa procedural} alrededor
del entrevistador-avatar reemplaza al panel de métricas tradicional y
materializa la metáfora del \textit{dojo}, sosteniendo la inmersión
durante la sesión (cubre \rf{05} y \rf{15}). Segundo, la totalidad del
análisis de video ocurre en el navegador mediante MediaPipe, sin envío
de imagen cruda a servicios externos, lo que mejora simultáneamente la
privacidad y el costo operativo (\rnf{05}). Tercero, la sesión
\textit{peer mock} sobre WebRTC habilita observación instrumentada
con métricas cruzadas y comentarios anclados a \textit{timestamps},
una capacidad ausente en Pramp y en los productos asistidos por IA
relevados (\rf{10}--\rf{12}). Cuarto, la accesibilidad se trata como
requerimiento de primer nivel y no como adición tardía: todos los
canales del aura tienen sustituto auditivo o textual y la operación
total del simulador es posible con teclado, conmutadores o voz
(\rf{17}--\rf{20} y \rnf{08}).

Este posicionamiento orienta también el alcance del MVP: el equipo
prioriza completar las capacidades diferenciadoras antes de replicar
funciones ya bien resueltas por la competencia (p.\,ej.\ catálogos
extensos de preguntas curados manualmente, formación en redacción de
\textit{currícula}, integración con plataformas de empleabilidad
externas), las cuales quedan documentadas como extensiones futuras.

\section{Criterios concretos de usabilidad}
\label{sec:usabilidad}

Esta sección consolida los criterios de usabilidad que orientarán el
diseño y la evaluación del simulador. La consolidación se hace en tres
capas complementarias: (i) los cinco atributos clásicos de
\textit{usability} de Nielsen (aprendizaje, eficiencia, memorabilidad,
manejo de errores, satisfacción), reformulados para el dominio de la
entrevista por voz; (ii) un subconjunto de las heurísticas de Nielsen
de evaluación heurística (HE), aplicadas a las pantallas de la sala
virtual y el plan de mejora; y (iii) un conjunto de \textbf{indicadores
verificables} (umbrales numéricos o conformidades normativas) que se
medirán durante el desarrollo y las pruebas con usuarios. Cada
indicador se enlaza con un requerimiento de la sección
\ref{sec:requerimientos} para asegurar trazabilidad.

\subsection{Atributos de usabilidad y traducción al dominio}

\begin{itemize}
  \item \textbf{Aprendizaje (\textit{learnability}).} Un usuario nuevo
        debe poder iniciar su primera sesión completa en menos de dos
        minutos desde el alta, sin recurrir al tutorial guiado. La
        metáfora del \textit{dojo} y la línea de tiempo circular se
        eligen explícitamente para acelerar la formación del modelo
        mental. Indicador: \rnf{10} \textit{(time-to-first-session)}.

  \item \textbf{Eficiencia (\textit{efficiency}).} Tras tres sesiones
        exitosas, el usuario debe poder iniciar una nueva entrevista en
        no más de dos clics o un comando de voz, gracias a las
        plantillas por industria persistidas en su perfil. Las acciones
        durante la sesión están disponibles tanto por puntero como por
        comando de voz reservado, lo que reduce la fricción al gesto
        físico. Indicador: \rf{06}, \rf{19}.

  \item \textbf{Memorabilidad (\textit{memorability}).} La disposición
        espacial del avatar y el aura es invariante entre los modos
        individual, peer mock y panel; lo único que cambia es la
        información presentada. Los \textit{badges} y rangos por
        competencia refuerzan el aprendizaje longitudinal. Indicador:
        \rf{13}, satisfacción reportada en pruebas con usuarios.

  \item \textbf{Manejo de errores (\textit{error handling}).} La
        interfaz debe prevenir errores costosos antes de cometerlos
        (confirmación explícita antes de publicar al \textit{feed},
        botón ``rehacer última respuesta''), recuperarlos sin pérdida
        de progreso (cola de reintentos ante caídas del LLM,
        reconexión TURN ante caída de WebRTC) y comunicarlos en
        lenguaje natural por voz del entrevistador, no como diálogos
        modales técnicos. Indicador: \rnf{04}, \rnf{14}.

  \item \textbf{Satisfacción.} La experiencia debe celebrar el proceso
        formativo y no únicamente el resultado, evitando los
        \textit{dark patterns} de comparación absoluta. La gamificación
        prioriza mejora relativa del usuario respecto a su línea base,
        no \textit{rankings} de pares ($\rf{14}$). La transparencia
        sobre el funcionamiento del análisis automático y sus
        eventuales sesgos ($\rnf{14}$) sostiene la confianza
        longitudinal en la herramienta.
\end{itemize}

\subsection{Heurísticas de Nielsen aplicadas al simulador}

Se aplican las diez heurísticas de Nielsen como rejilla de evaluación
durante la revisión interna del prototipo y como criterios formales
para la auditoría heurística externa antes de la primera prueba con
usuarios. La tabla \ref{tab:heuristicas} indica, para cada heurística,
cómo se materializa en el simulador y en qué pantalla se verifica.

\begin{table}[ht]
\centering
\renewcommand{\arraystretch}{1.25}
\footnotesize
\begin{tabularx}{\linewidth}{@{}p{4.5cm} Y l@{}}
\toprule
\textbf{Heurística} & \textbf{Materialización en el simulador} &
\textbf{Pantalla / RF} \\
\midrule
H1. Visibilidad del estado del sistema &
El aura procedural y la línea de tiempo circular comunican el estado
de la sesión y del análisis sin abrir paneles separados. &
Sala virtual / \rf{05} \\
\addlinespace
H2. Correspondencia entre sistema y mundo real &
Lenguaje neutral propio del proceso de selección laboral; metáfora
del \textit{dojo} familiar para el usuario. &
Sala virtual, plan de mejora \\
\addlinespace
H3. Control y libertad del usuario &
Comandos ``pausa'', ``repetir'', ``terminar'' siempre disponibles;
botón ``rehacer última respuesta'' antes del cierre. &
Sala virtual / \rf{06} \\
\addlinespace
H4. Consistencia y estándares &
Misma disposición espacial en los modos individual, peer y panel;
componentes de aura y \textit{tooltip} reutilizables. &
Toda la aplicación \\
\addlinespace
H5. Prevención de errores &
Confirmación explícita antes de publicar grabaciones al \textit{feed};
\textit{onboarding} de permisos antes de la primera sesión. &
Privacidad, comunidad / \rf{21} \\
\addlinespace
H6. Reconocer mejor que recordar &
\textit{Tooltips} contextuales sobre cada icono del aura con criterio
y nivel de confianza; sugerencias del asistente lateral basadas en
historial. &
Sala virtual, asistente / \rf{15}, \rf{16} \\
\addlinespace
H7. Flexibilidad y eficiencia &
Atajos de voz y de teclado declarados en cada control; plantillas
por industria; navegación 100\,\% por teclado para todo flujo. &
Configuración / \rf{19} \\
\addlinespace
H8. Diseño estético y minimalista &
La sala virtual oculta cualquier información no esencial durante la
sesión; las métricas viven en el aura, no en \textit{dashboards}. &
Sala virtual / \rf{05} \\
\addlinespace
H9. Recuperación de errores en lenguaje claro &
Mensajes del entrevistador en español neutro, sin códigos técnicos;
explicación verbal de caídas de red o del LLM. &
Sala virtual / \rnf{14} \\
\addlinespace
H10. Ayuda y documentación &
Tour guiado al primer ingreso, asistente lateral por demanda,
ayuda inferida por inactividad o reintentos repetidos. &
Toda la aplicación / \rf{15}, \rf{16} \\
\bottomrule
\end{tabularx}
\caption{Aplicación de las diez heurísticas de Nielsen al simulador.}
\label{tab:heuristicas}
\end{table}

\subsection{Indicadores verificables}

Los criterios anteriores se traducen en indicadores con umbrales o
conformidades concretas, evaluables durante el desarrollo (mediante
prueba automatizada o telemetría) y durante las pruebas con usuarios
(mediante observación cronometrada y cuestionarios SUS). La tabla
\ref{tab:indicadores} resume el conjunto.

\begin{table}[ht]
\centering
\renewcommand{\arraystretch}{1.25}
\footnotesize
\begin{tabularx}{\linewidth}{@{}p{3.6cm} Y p{3.0cm} p{2.4cm}@{}}
\toprule
\textbf{Indicador} & \textbf{Definición operativa} &
\textbf{Umbral / objetivo} & \textbf{Trazabilidad} \\
\midrule
Tiempo a primera sesión & Tiempo entre alta del usuario y final de su
primera entrevista. & $\le 2$ minutos & \rnf{10} \\
\addlinespace
Latencia del aura & Retardo extremo a extremo entre evento percibido y
actualización visual del aura. & $\le 500$\,ms (p95) & \rnf{01},
\rf{05} \\
\addlinespace
Latencia de respuesta del LLM & Retardo entre fin de turno del
candidato y comienzo de la voz del entrevistador. & $\le 2$\,s (p95)
& \rnf{02}, \rf{02} \\
\addlinespace
Latencia de subtítulos & Retardo de aparición del subtítulo respecto
al fonema de inicio de la pregunta del LLM. & $\le 300$\,ms (p95) &
\rf{18} \\
\addlinespace
Tasa de finalización de tarea & Porcentaje de usuarios que completan
una entrevista de extremo a extremo sin asistencia humana. &
$\ge 90\%$ tras 30\,s de inducción & Pruebas moderadas \\
\addlinespace
Tasa de error en STT &
Word Error Rate sobre transcripciones del candidato en español neutro.
& $\le 12\%$ (referencial) & \rf{03} \\
\addlinespace
Conformidad WCAG 2.2 AA &
Auditoría \texttt{axe-core} en CI sin violaciones \textit{serious} ni
\textit{critical}. & 0 violaciones & \rnf{08} \\
\addlinespace
Puntaje Lighthouse &
Categoría \textit{Accessibility} ejecutada en CI sobre las cinco rutas
críticas. & $\ge 95$ & \rnf{08} \\
\addlinespace
Conformidad de teclado &
Cobertura E2E \textit{keyboard-only} de los flujos de usuario. &
$100\%$ de flujos críticos & \rf{19}, \rnf{08} \\
\addlinespace
SUS (System Usability Scale) &
Cuestionario aplicado a una muestra de candidatos primarios al cierre
de la prueba con usuarios. & $\ge 75$ & Prueba final con
usuarios \\
\addlinespace
Costo de IA por sesión &
Costo medio en USD por sesión completa de 30 minutos. &
$\le \$0{,}20$ & \rnf{13} \\
\bottomrule
\end{tabularx}
\caption{Indicadores verificables de usabilidad y conformidad,
con umbrales operativos y trazabilidad a requerimientos.}
\label{tab:indicadores}
\end{table}

\subsection{Estrategia de evaluación}

La verificación de los indicadores anteriores se distribuye en cuatro
hitos a lo largo del ciclo de desarrollo: (i) revisión heurística
interna por dos miembros del equipo siguiendo la rejilla de la tabla
\ref{tab:heuristicas} antes de la primera prueba con usuarios,
(ii) ejecución continua de \texttt{axe-core} y \textit{Lighthouse} en
\textit{integración continua} sobre cada \textit{pull request} para
los indicadores de accesibilidad y rendimiento, (iii) pruebas
moderadas con un usuario representante por perfil primario
(\sk{1}--\sk{7}) sobre una entrevista completa, registrando los
indicadores conductuales (tasa de finalización, tiempo a primera
sesión) y aplicando el cuestionario SUS al cierre, y
(iv) auditoría externa de accesibilidad previa a la entrega final, con
informe firmado de conformidad WCAG 2.2 nivel AA.

Los criterios anteriores no son una capa decorativa: cada uno se
declara como entrada al \textit{backlog} de pruebas automatizadas o al
plan de pruebas con usuarios, de modo que un fallo medible en
cualquiera de ellos bloquea la aceptación del módulo correspondiente.


\end{document}
